%%%PAGE 222 rot=0 legibility=good summary=解题方法·参考系与相对运动（相对速度、槽与小球、错车时间、管球自由落体、向心力用相对圆心速度、E=BL(V1-V2)）
%%%BLOCK id=222-01 ch=01 sec=参考系·相对运动 kind=knowledge alt=none cont=none y=0.04-0.18
\textbf{2、参考系}

运动的相对性\quad $V_{\text{相}}=\vec V_{\text{甲乙}}=\vec V_{\text{甲}}-\vec V_{\text{乙}}$

选参的任意性、同一性\quad 先方便后地面。

定正方向用±号表示方向
%%%END
%%%BLOCK id=222-02 ch=07 sec=动量与能量综合·曲面槽模型 kind=method alt=06,01 cont=none y=0.18-0.29
%FIG p222_a 0.08 0.18 0.395 0.285
\notefig[0.4]{figures/p222_a.png}
$V_{\text{临}}=\sqrt{\dfrac{M+m}{M}}\sqrt{2gR}$ % CHECK: 根号内分子被涂改，难辨（M+m 或 m+M）

仍会落在槽内

$V_{x\text{相}}=0$。若到最高点时将槽固定，

$V_{x\text{相}}=V_{\text{共}}$
%%%END
%%%BLOCK id=222-03 ch=01 sec=相对运动·错车时间 kind=problem alt=none cont=none y=0.29-0.43
%FIG p222_b 0.08 0.295 0.395 0.38
\notefig[0.4]{figures/p222_b.png}
错开的时间
\[
t=\frac{x_{\text{甲乙}}}{V_{\text{甲乙}}}=\frac{14}{5}=2.8\unit{s}
\] % CHECK: 分母首字被涂改（t 改 V？）
%%%END
%%%BLOCK id=222-04 ch=01 sec=参考系·相对运动 kind=knowledge alt=none cont=none y=0.44-0.49
若 $a_1=a_2$，则 1、和 2 相对静止或相对匀速
%%%END
%%%BLOCK id=222-05 ch=01 sec=自由落体·相对运动 kind=method alt=none cont=none y=0.49-0.63
%FIG p222_c 0.055 0.505 0.22 0.612
\notefig[0.25]{figures/p222_c.png}
管在球到\unclear{空}时放手

$t_{\text{管中}}=\dfrac{x}{v}=\unclear{1}\unit{s}$
%%%END
%%%BLOCK id=222-06 ch=04 sec=圆周运动·向心力（相对速度） kind=method alt=07,06 cont=none y=0.64-0.83
向心力。相对于圆心的速度\quad $F_n=\dfrac{mv^2}{r}$

%FIG p222_d 0.0 0.705 0.15 0.785
\notefig[0.25]{figures/p222_d.png}
\begin{align*}
mgR&=\frac12MV'^2+\frac12mV^2\\
mV&=MV'\\
N-mg&=\frac{m(V+V')^2}{r}
\end{align*}
%%%END
%%%BLOCK id=222-07 ch=12 sec=动生电动势·相对速度 kind=method alt=none cont=none y=0.83-0.91
%FIG p222_e 0.005 0.83 0.205 0.90
\notefig[0.25]{figures/p222_e.png}
$E=BL(V_1-V_2)$\quad 导体棒相对于磁场的速度。
%%%END
%%%PAGEEND 222
